\documentclass[11pt]{article}
\usepackage[margin=1in]{geometry}
\usepackage{amsmath,amssymb,amsthm}
\usepackage{hyperref}
\newtheorem{theorem}{Theorem}
\newtheorem{lemma}[theorem]{Lemma}
\newtheorem{proposition}[theorem]{Proposition}
\theoremstyle{definition}
\newtheorem{definition}[theorem]{Definition}
\theoremstyle{remark}
\newtheorem{remark}[theorem]{Remark}
\newcommand{\F}{\mathbb{F}_8}
\newcommand{\pol}{f}

\title{Disproof of Conjecture 00000001007:\\
the Suzuki--Tits ovoid of $Q(4,8)$ is not an elliptic quadric}
\author{jilint777}
\date{2026-10-04}

\begin{document}
\sloppy
\maketitle

\begin{abstract}
Conjecture 00000001007 asserts that in even characteristic every ovoid of the
parabolic quadric $Q(4,q)$ is an elliptic quadric. This is false for $q=8$.
The Suzuki--Tits ovoid (Tits, 1962), lifted to $Q(4,8)$, is an ovoid: it has
$q^2+1=65$ points and meets every line of $Q(4,8)$ exactly once. But it lies in
no hyperplane of $\mathrm{PG}(4,8)$, so it is not a hyperplane section
$Q^-(3,8)$ of $Q(4,8)$. Its projection from the nucleus to $\mathrm{PG}(3,8)$
lies on no quadric at all, so it is not an elliptic quadric in that model
either. All facts are checked in Lean~4 (core library only) and by an
independent Python script. Short hand proofs are also given.
\end{abstract}

\section{The conjecture and what we refute}
The statement reads:
\begin{quote}
\emph{Conjecture:} In even characteristic, every ovoid of $Q(4,q)$ is an
elliptic quadric (uniqueness; the odd-characteristic case has Suzuki's
counterexamples) --- the even-characteristic uniqueness conjecture.
\end{quote}
The Chinese version says the same thing. The clause we refute is
\begin{center}
(C)\quad \emph{for every even prime power $q$, every ovoid of $Q(4,q)$ is an
elliptic quadric.}
\end{center}
One even value of $q$ with a non-elliptic ovoid is enough, and we use $q=8$.
The parenthetical remark is not a mathematical claim we need. It is also
inaccurate: the Suzuki--Tits ovoids exist only in characteristic $2$, for
$q=2^{2e+1}\ge 8$. They give counterexamples to (C). (Whether they are the only non-classical
ovoids of $\mathrm{PG}(3,q)$, $q$ even, is open in general.)

\paragraph{Readings.}
``Ovoid of $Q(4,q)$'' has a standard meaning (Definition~\ref{def:ovoid}). We
also use the standard equivalent finite form (Lemma~\ref{lem:equiv}).
``Elliptic quadric'' can be read in two models.
\begin{itemize}
\item[(R1)] \emph{Inside $\mathrm{PG}(4,q)$.} An elliptic quadric ovoid of
$Q(4,q)$ is a section $Q(4,q)\cap H$ by a hyperplane $H$ such that $Q|_H$ is
elliptic ($Q^-(3,q)$). The same holds for any image of such a section under a
collineation or a semilinear isometry, since these maps send hyperplanes to
hyperplanes. Every such set lies in a hyperplane of $\mathrm{PG}(4,q)$.
\item[(R2)] \emph{Inside $\mathrm{PG}(3,q)$.} For $q$ even, projection from the
nucleus $N$ of $Q(4,q)$ is a bijection from the points of $Q(4,q)$ onto the
points of $\mathrm{PG}(3,q)$, and ovoids of $Q(4,q)$ map to ovoids of
$\mathrm{PG}(3,q)$. In this model ``elliptic quadric'' means the zero set of an
elliptic quadratic form on $\F^4$, possibly after a collineation. Every such
set lies on a quadric, i.e.\ on the zero set of a nonzero quadratic form,
since collineations map quadrics to quadrics.
\end{itemize}
We exhibit an ovoid $\mathcal{T}$ of $Q(4,8)$ that lies in no hyperplane and
whose projection lies on no quadric. This defeats both readings and any reading
that implies one of these two necessary conditions.
We refute the statement as literally written, in both the English and the
Chinese versions. A different statement, such as ``every ovoid other than the
Suzuki--Tits ovoids is elliptic'', is the open classification problem, and we
make no claim about it. The literal parenthetical places Suzuki's examples in
odd characteristic, so it gives no exception for even $q$.

\section{The objects}
\paragraph{The field.}
$\F=\mathbb{F}_2[\omega]/(\omega^3+\omega+1)$. We code $b_0+b_1\omega+b_2\omega^2$
by the integer $b_0+2b_1+4b_2$. So $2=\omega$, $4=\omega^2$,
$3=\omega^3$, $6=\omega^4$, $7=\omega^5$, $5=\omega^6$, and $\omega$ has
order $7$. Addition is XOR of codes, and $\operatorname{char}\F=2$, so $q=8$ is
even. The map $\sigma(x)=x^4$ is a field automorphism with
$\sigma(\sigma(x))=x^{16}=x^2$, so $\sigma^2$ is the Frobenius map. This is the
automorphism needed for the Tits ovoid when $q=2^{2e+1}$ with $e=1$.

\paragraph{The quadric.}
On $V=\F^5$ let
\[
Q(x)=x_0x_3+x_1x_2+x_4^2,\qquad
\pol(x,y)=x_0y_3+x_3y_0+x_1y_2+x_2y_1 .
\]
The identities $(a+b)(c+d)=ac+bd+(ad+bc)$ and $(a+b)^2=a^2+b^2$ give
$Q(x+y)=Q(x)+Q(y)+\pol(x,y)$, so $\pol$ is the polar form of $Q$. Also
$\pol(x,x)=2(x_0x_3+x_1x_2)=0$. Pairing with the unit vectors
$e_3,e_2,e_1,e_0$ gives $\pol(x,\cdot)=0$ if and only if
$x_0=x_1=x_2=x_3=0$. So the radical of $\pol$ is spanned by the nucleus
$N=(0,0,0,0,1)$, and $Q(N)=1\ne0$. Hence $Q$ is a nondegenerate quadratic form
of parabolic type, and its zero set $Q(4,8)$ in $\mathrm{PG}(4,8)$ is the
parabolic quadric. All nondegenerate quadrics of $\mathrm{PG}(4,q)$ are
projectively equivalent, so the choice of model loses nothing: an isometry maps
ovoids to ovoids and hyperplane sections to hyperplane sections.

For points $x\ne y$ of $Q(4,8)$,
$Q(sx+ty)=s^2Q(x)+t^2Q(y)+st\,\pol(x,y)=st\,\pol(x,y)$. So the line $xy$ lies on
the quadric if and only if $\pol(x,y)=0$. Such points are called
\emph{collinear}. The lines of $Q(4,8)$ are the totally singular
$2$-dimensional subspaces of $V$.

\paragraph{Counts.}
$Q(4,q)$ is a generalized quadrangle of order $(q,q)$ \cite[\S3.1]{PT}. It has
$(q+1)(q^2+1)$ points and $(q+1)(q^2+1)$ lines. Each line has $q+1$ points and
each point lies on $q+1$ lines. For $q=8$ there are $585$ points, $585$ lines,
$9$ points per line and $9$ lines per point. \texttt{verify.py} confirms these
numbers by listing all $4681$ points of $\mathrm{PG}(4,8)$ and all lines of the
quadric.

\begin{definition}\label{def:ovoid}
An \emph{ovoid} of $Q(4,q)$ is a set $\mathcal{O}$ of points of $Q(4,q)$ that
meets every line of $Q(4,q)$ in exactly one point.
\end{definition}

\begin{lemma}\label{lem:equiv}
A set $\mathcal{O}$ of points of $Q(4,q)$ is an ovoid if and only if
$|\mathcal{O}|=q^2+1$ and no two points of $\mathcal{O}$ are collinear.
\end{lemma}
\begin{proof}
($\Leftarrow$) Two points of $\mathcal{O}$ on a common line of $Q(4,q)$ would
be collinear, so each line contains at most one point of $\mathcal{O}$. Each
point of $\mathcal{O}$ lies on $q+1$ lines, and these lines are distinct for
distinct points of $\mathcal{O}$. So $\mathcal{O}$ meets at least
$(q^2+1)(q+1)$ lines, which is all of them.
($\Rightarrow$) Two collinear points of $\mathcal{O}$ would put two points of
$\mathcal{O}$ on one line. Counting incidences between $\mathcal{O}$ and lines
gives $|\mathcal{O}|(q+1)=(q+1)(q^2+1)$.
\end{proof}
Lean formalizes the right-hand side as \texttt{IsOvoid}, for $q=8$. It also
proves the ``at most one point per line'' half of ($\Leftarrow$) for every
totally singular line (\texttt{IsOvoid.at\_most\_one}). The counting half is
the argument above. \texttt{verify.py} checks the defining property directly:
each of the $585$ lines meets our set in exactly one point.

\section{The Suzuki--Tits ovoid}
In $\mathrm{PG}(3,8)$ the Tits ovoid \cite{Tits} is
\[
\mathcal{T}_3=\{(1,x,y,\;xy+x^{\sigma+2}+y^{\sigma}) : x,y\in\F\}\cup\{(0,0,0,1)\},
\qquad x^{\sigma+2}=x^6,\ y^\sigma=y^4 .
\]
We lift it to $Q(4,8)$. Every $a\in\F$ has a unique square root, because
squaring is an automorphism. We append to $(p_0,p_1,p_2,p_3)$ the coordinate
$\sqrt{p_0p_3+p_1p_2}$. For an affine point this square root is
$\sqrt{x^6+y^4}=x^3+y^2$. So
\[
\mathcal{T}=\{t(x,y):=(1,\,x,\,y,\,xy+x^6+y^4,\,x^3+y^2) : x,y\in\F\}\cup\{(0,0,0,1,0)\}.
\]
The point $t(x,y)$ is on the quadric because
$Q(t(x,y))=(xy+x^6+y^4)+xy+(x^3+y^2)^2=0$. The last point is on the quadric
because $Q(0,0,0,1,0)=0$. Projection from $N$ (deleting $x_4$) maps
$\mathcal{T}$ onto $\mathcal{T}_3$.

\begin{proposition}\label{prop:ovoid}
$\mathcal{T}$ is an ovoid of $Q(4,8)$. It consists of $65$ distinct points, no
two of them collinear.
\end{proposition}
\begin{proof}
The $64$ points $t(x,y)$ have distinct coordinates $(x_1,x_2)$ and first
coordinate $1$. The point $(0,0,0,1,0)$ is the only one with $x_0=0$. So the
65 points are pairwise non-proportional. Moreover
$\pol\big((0,0,0,1,0),t(x,y)\big)=1$, and for $(x,y)\ne(x',y')$,
\[
\pol\big(t(x,y),t(x',y')\big)=(x+x')(y+y')+x^6+x'^6+(y+y')^4 .
\]
That this is nonzero is Tits' theorem \cite{Tits}. Here it is checked
exhaustively for all $2080$ pairs, by \texttt{decide} in Lean and in
\texttt{verify.py}. By Lemma~\ref{lem:equiv}, $\mathcal{T}$ is an ovoid.
\texttt{verify.py} also checks Definition~\ref{def:ovoid} directly on all
585 lines, and checks that $\mathcal{T}_3$ has no three collinear points (all
$43680$ triples), so $\mathcal{T}_3$ is an ovoid of $\mathrm{PG}(3,8)$.
\end{proof}

\section{\texorpdfstring{$\mathcal{T}$ is not an elliptic quadric}{T is not an elliptic quadric}}
We use one elementary fact: a polynomial over $\F$ of degree at most $7$ that
vanishes at all $8$ elements of $\F$ is the zero polynomial.

\begin{proposition}[reading R1]\label{prop:hyp}
No hyperplane of $\mathrm{PG}(4,8)$ contains $\mathcal{T}$.
\end{proposition}
\begin{proof}
Let $a\in\F^5$ satisfy $\sum_i a_ip_i=0$ for all $p\in\mathcal{T}$. The points
$t(0,y)=(1,0,y,y^4,y^2)$ give $a_0+a_2y+a_4y^2+a_3y^4=0$ for all $y\in\F$. This
polynomial has degree at most $4$, so $a_0=a_2=a_3=a_4=0$. Then
$t(1,0)=(1,1,0,1,1)$ gives $a_1=0$. So $a=0$.
\end{proof}

\begin{proposition}[reading R2]\label{prop:quad}
No nonzero quadratic form $c(x)=\sum_{0\le i\le j\le 3}c_{ij}x_ix_j$ on $\F^4$
vanishes on $\mathcal{T}_3$. Hence $\mathcal{T}_3$ lies on no quadric of
$\mathrm{PG}(3,8)$.
\end{proposition}
\begin{proof}
The point $(0,0,0,1)$ gives $c_{33}=0$. The points $(1,0,y,y^4)$ give
$c_{00}+c_{02}y+c_{22}y^2+c_{03}y^4+c_{23}y^5=0$ for all $y$. This polynomial has
degree at most $5$, so these five coefficients vanish. The points
$(1,x,0,x^6)$ now give $c_{01}x+c_{11}x^2+c_{13}x^7=0$ for all $x$, since
the terms with $c_{00},c_{03},c_{33}$ are already zero. This polynomial has
degree at most $7$, so $c_{01}=c_{11}=c_{13}=0$. Finally the point
$(1,1,1,1)=t(1,1)$ with the last coordinate deleted gives $c_{12}=0$.
\end{proof}

\begin{theorem}
Clause (C) is false: $\mathcal{T}$ is an ovoid of $Q(4,8)$, $q=8$ is even, and
$\mathcal{T}$ is not an elliptic quadric under reading R1 or R2.
\end{theorem}
\begin{proof}
Proposition~\ref{prop:ovoid} shows that $\mathcal{T}$ is an ovoid. An elliptic
quadric under R1 lies in a hyperplane, which Proposition~\ref{prop:hyp} rules
out. Under R2 its projection lies on a quadric, which
Proposition~\ref{prop:quad} rules out.
\end{proof}

\paragraph{Control (the predicates are not vacuous).}
The polynomial $x^2+x+1$ has no root in $\F$. So the hyperplane
$H:\,x_4=x_1+x_2$ cuts $Q(4,8)$ in the elliptic quadric
$x_0x_3+x_1^2+x_1x_2+x_2^2=0$. That section is
$\mathcal{E}=\{(1,a,b,a^2+ab+b^2,a+b)\}\cup\{(0,0,0,1,0)\}$. Lean proves that
$\mathcal{E}$ satisfies \texttt{IsOvoid}, lies in $H$, and lies on the quadric
$x_0x_3+x_1^2+x_1x_2+x_2^2$ of $\mathrm{PG}(3,8)$. \texttt{verify.py} checks
that $\mathcal{E}=Q(4,8)\cap H$ and that it meets every line once.
$\mathcal{T}$ and $\mathcal{E}$ share only $5$ points.

\begin{remark}
For $q=2$ and $q=4$ every ovoid of $\mathrm{PG}(3,q)$ is an elliptic quadric
(classical results of Barlotti and others), and hence so is every ovoid of $Q(4,q)$. So $q=8$ is the smallest
counterexample. For every $q=2^{2e+1}\ge8$ the Tits ovoid gives a counterexample
in the same way. For odd prime $q$, every ovoid of $Q(4,q)$ is an elliptic quadric
(Ball, Govaerts and Storme). None of this is used above.
\end{remark}

\section{Lean formalization}
The project \texttt{lean4/} uses Lean 4.19.0 and the core library only. It
contains no \texttt{sorry}, no \texttt{native\_decide} and no added axioms.
\begin{itemize}
\item \texttt{F8}: the eight elements, with addition and multiplication tables
generated from $\omega^3=\omega+1$. The theorems \texttt{add\_assoc},
\texttt{mul\_inv}, \texttt{left\_distrib} and so on check the field axioms
exhaustively. \texttt{char\_two} shows that the characteristic is $2$, and
\texttt{omega\_generates} shows that $\omega$ has order $7$.
\texttt{sigma\_props} shows that $\sigma(x)=x^4$ is an automorphism with
$\sigma^2=$ Frobenius.
\item \texttt{Q} and \texttt{pol} are $Q$ and $\pol$. The theorem
\texttt{Q\_vadd} is the polarization identity, proved by algebra.
\texttt{Q\_nondegenerate} shows that the radical is
$\langle N\rangle$, that $Q(N)=1$, and that no nonzero singular vector lies in
the radical.
\item \texttt{IsOvoid O}: $O$ has $65$ distinct entries, all nonzero and on $Q$,
and $\pol(p,p')\ne0$ for distinct entries. \texttt{IsOvoid.proj\_distinct}
shows that the entries are distinct projective points.
\texttt{IsOvoid.at\_most\_one} shows that every totally singular line
$\{sx+ty\}$ contains at most one entry.
\item \texttt{tits} is $\mathcal{T}$, written with \texttt{sigma}.
\texttt{tits\_lift} shows that $(x^3+y^2)^2=x^{\sigma+2}+y^\sigma$, and
\texttt{tits\_isOvoid} is proved by \texttt{decide +kernel}.
\item \texttt{InHyperplane O} and \texttt{OnQuadricPG3 O} are the necessary
conditions of R1 and R2. \texttt{tits\_not\_inHyperplane} and
\texttt{tits\_not\_onQuadric} prove that $\mathcal{T}$ fails both. These proofs
do not enumerate the $8^5$ or $8^{10}$ coefficient vectors. They use a general
lemma (\texttt{dotL\_comb}, \texttt{kill}): if $c$ is orthogonal to the rows of
$E$ and $NE=I$, then $c=0$. They also use explicit inverse matrices $N$ for $5$
points (linear forms) and $10$ points (quadratic monomials) of $\mathcal{T}$,
checked by \texttt{decide}.
\item \texttt{conjecture\_00000001007\_false}:
\texttt{$\neg$(\,$\forall$ O, IsOvoid O $\to$ InHyperplane O)}.
\texttt{conjecture\_00000001007\_false\_any\_reading}: for every predicate
\texttt{IsElliptic} whose instances satisfy \texttt{InHyperplane} or
\texttt{OnQuadricPG3}, \texttt{$\neg$(\,$\forall$ O, IsOvoid O $\to$ IsElliptic O)}.
\texttt{nonvacuous} and \texttt{ell\_isOvoid} show that $\mathcal{E}$ satisfies
all of these predicates.
\end{itemize}
\paragraph{Limitations.}
Lean fixes $q=8$ and the standard model of $Q(4,8)$; it does not formalize
general finite fields. This is enough, because (C) quantifies over all even
$q$. The ovoid predicate in Lean is the finite form of Lemma~\ref{lem:equiv}.
Lean proves the ``at most one point per line'' half of that lemma. The
counting half and the generalized-quadrangle counts are proved above, and
\texttt{verify.py} checks the line definition exhaustively. ``Elliptic
quadric'' enters Lean only through its necessary conditions. Refuting those
weaker conditions gives a stronger disproof.

\section{Reproduction}
\begin{verbatim}
cd lean4 && lake build        # about 40 s
cd .. && python3 verify.py    # about 10 s, standard library only
\end{verbatim}

\begin{thebibliography}{9}
\bibitem{Tits} J.~Tits, \emph{Ovo\"ides et groupes de Suzuki},
Arch.\ Math.\ 13 (1962), 187--198.
\bibitem{PT} S.~E.~Payne and J.~A.~Thas, \emph{Finite Generalized Quadrangles},
2nd ed., EMS, 2009.
\end{thebibliography}
\end{document}
